% !TEX root = task-2-10-4.tex
% Editable task bank. Append new BankTask blocks at the end of the document.
\documentclass[12pt,a4paper]{article}
\usepackage{physics-textbook}
\usepackage{graphicx,needspace}
\usetikzlibrary{decorations.pathmorphing}
\newif\ifShowSolutions
\ifdefined\TutorEdition
  \ShowSolutionstrue
\else
  \ShowSolutionsfalse % Change to \ShowSolutionstrue to display solutions.
\fi
\newcounter{banktask}
\newsavebox{\BankTaskBox}
\newcommand{\BankTaskContent}[3]{%
  \noindent\textbf{\thebanktask.}\quad #2\par
  \ifShowSolutions
    \Needspace{50mm}% Keep the solution heading with its opening text or diagram.
    #3\par
    \ifdefined\TutorEdition
      \ifcsname TutorSource#1\endcsname
        \nopagebreak\smallskip{\small\csname TutorSource#1\endcsname\par}%
      \fi
    \fi
  \fi
}
\newcommand{\BankTask}[3]{%
  \refstepcounter{banktask}\par\medskip
  \sbox{\BankTaskBox}{\begin{minipage}{\linewidth}%
    \BankTaskContent{#1}{#2}{#3}\end{minipage}}%
  \ifdim\dimexpr\ht\BankTaskBox+\dp\BankTaskBox\relax>\dimexpr\textheight-20mm\relax
    \Needspace{35mm}\BankTaskContent{#1}{#2}{#3}%
  \else
    \Needspace{\dimexpr\ht\BankTaskBox+\dp\BankTaskBox\relax}%
    \noindent\usebox{\BankTaskBox}\par
  \fi
  \medskip
}
\begin{document}
\BookSubsection{2.10.4}{График кристаллизации}


% Source 35852E: https://ege.fipi.ru/bank/questions.php?search=1&proj=BA1F39653304A5B041B656915DC36B38&qid=35852E
\BankTask{35852E}{
Горячее вещество, первоначально находившееся в жидком состоянии, медленно охлаждали. Мощность теплоотвода постоянна. В таблице приведены результаты измерений температуры вещества с течением времени.
\begin{center}
\begin{tabular}{|l|r|r|r|r|r|r|r|r|}
\hline
Время, $\text{мин}$ & \(0\) & \(5\) & \(10\) & \(15\) & \(20\) & \(25\) & \(30\) & \(35\) \\
\hline
Температура, \({}^{\circ}\mathrm{C}\) & \(250\) & \(242\) & \(234\) & \(232\) & \(232\) & \(232\) & \(230\) & \(216\) \\
\hline
\end{tabular}
\end{center}
\begin{enumerate}
\item Температура плавления вещества в данных условиях равна \(232\,{}^{\circ}\mathrm{C}\).
\item Через \(20\,\text{мин}\). после начала измерений вещество находилось только в твёрдом состоянии.
\item Удельная теплоёмкость вещества в жидком и твёрдом состояниях одинакова.
\item Через \(30\,\text{мин}\). после начала измерений вещество находилось только в твёрдом состоянии.
\item Процесс кристаллизации вещества занял более \(25\,\text{мин}\).
\end{enumerate}
}{
\smallskip\noindent\textbf{Решение от GPT}\par
\begin{enumerate}
\item Верно: участок постоянной температуры соответствует кристаллизации, поэтому \(T_{\text{пл}}=232\,{}^{\circ}\mathrm{C}\).
\item Неверно: в момент \(t=20\,\text{мин}\) идёт кристаллизация, присутствуют жидкая и твёрдая фазы.
\item Неверно: при постоянной мощности \(P=mc\lvert\Delta T/\Delta t\rvert\). Скорости охлаждения жидкости и твёрдого вещества равны соответственно \(8/5=1{,}6\) и \(14/5=2{,}8\,{}^{\circ}\mathrm{C}/\text{мин}\); теплоёмкости различны.
\item Верно: через \(30\,\text{мин}\) температура \(230\,{}^{\circ}\mathrm{C}\) ниже температуры кристаллизации, вещество полностью затвердело.
\item Неверно: кристаллизация началась после \(10\,\text{мин}\) и завершилась до \(30\,\text{мин}\), поэтому её длительность меньше \(20\,\text{мин}\).
\end{enumerate}
Ответ: \(14\).
}

\end{document}
